\scanpage{296}
\begin{center}- 100 -\end{center}\hfill IX

9. \uline{Définition de l'accouplement de monodromie} $\underline M_\ell\otimes\underline M'_\ell\longrightarrow\mathbf Z_{\ell S}$.

9.1. Dans le présent numéro on suppose $S$ hensélien, et que $A_K$ ait réduction semi-stable sur $S$, ou ce qui revient au même, que $A'_K$ ait réduction semi-stable sur $S$ (3.5.1). On dispose alors sur $S$ de deux $\mathbf Z$-Modules constants tordus $\underline M$ et $M'$ (5.5.2), qui sur un revêtement fini étale convenable de $S$ (donc sur $S$ lui-même si $S$ est strictement local i.e. si $k$ est séparablement clos) se déploient en des groupes constants de valeur $M$ resp. $M'$ isomorphes (non canoniquement) à $\mathbf Z^r$, $r$ étant la dimension commune des tores maximaux de $A_o$ et de $A'_o$. Nous allons considérer, pour tout nombre premier $\ell$, les faisceaux $\ell$-adiques correspondants
\[
\tag{9.1.1}\underline M_\ell=\underline M\otimes_{\mathbf Z}\mathbf Z_\ell\ ,\quad
\underline M'_\ell=\underline M'\otimes_{\mathbf Z}\mathbf Z_\ell\ ,
\]
dont la signification géométrique est explicitée par les formules (5.5.7) et (5.5.8). Dans le présent numéro, nous allons définir un accouplement canonique (dit ``\uline{accouplement de monodromie}'')
\[
\tag{9.1.2}u_\ell:\underline M_\ell\otimes\underline M'_\ell\longrightarrow\mathbf Z_\ell\ ,
\]
qui sera défini en termes des pro-$\ell$-groupes de Barsotti-Tate $T_\ell(A^o)$ et $T_\ell(A'^o)$, et même de la fibre générique $T_\ell(A_K)$ du premier de ces groupes (sous réserve, lorsque $\ell=p$, de disposer du théorème de Tate [33, th.4], i.e. pour l'instant qu'on soit alors dans le cas d'inégales caractéristiques). Lorsque $\ell\ne p$, il revient au même de dire que l'accouplement (9.1.2) est décrit en termes du $\pi$-module galoisien
\[
\tag{9.1.3}U=T_\ell(A_K)(\overline K)=T_\ell(A_{\overline K})
\]
\margin{$\unsure{\ni}\qquad\unsure{u}$}
\margin{\figure{\begin{tikzpicture}\draw (0,0)--(.23,-.25);\draw (0,-.7)--(.3,-.72);\end{tikzpicture}}}
\margin{\unsure{418}}
\margin{418}
\margin{297}
